\section{Circuit Émetteur}

\subsection{Choix des composantes et explications}
\begin{frame}
	\frametitle{Choix des composantes}
	\begin{columns}
		\begin{column}{.5\textwidth}
			\begin{table}
				\caption{Composantes de l'émetteur déterminées}
				\begin{tabular}{ll}
					\toprule
					Résistance & Valeur Choisie  \\
					\midrule\midrule
					$R_6$      & \SI{1500}{\ohm} \\
					$R_7$      & \SI{8200}{\ohm} \\
					\bottomrule
				\end{tabular}
			\end{table}
		\end{column}

		\begin{column}{.5\textwidth}
			\begin{figure}
				\includegraphics[width=.7\textwidth]{figure/screenshot-emitter.jpg}
				\caption{Circuit de l'émetteur}
				\label{fig:circuit-emitter}
			\end{figure}
		\end{column}
	\end{columns}
\end{frame}

\subsection{À la charge}
\begin{frame}
	\frametitle{Mises en équation de la charge}
	\begin{align}
		V_\mathrm{in}(0) & = 12; \quad
		V_C(0)             = -12; \quad
		V_D(0)             = 0.6; \quad
		V_R(0)             = \SI{23.4}{\V}                               \\
		V_\mathrm{in}(t) & = V_C(t) + V_D(t) + V_R(t) \Rightarrow
		V_\mathrm{in}^{\prime}(t) = \frac{1}{RC}V_R(t) + V_R^{\prime}(t) \\
		V_R(t)           & = 23.4\e^{\frac{-t}{1.5\times10^{-5}}}
	\end{align}
\end{frame}

\begin{frame}
	\frametitle{Graphiques de charge théoriques}
	\begin{columns}

		\begin{column}{0.5\textwidth}
			\begin{figure}
				\includegraphics[width=.9\textwidth]{figure/dodo.jpg}
				\caption{M.E.É. de l'émetteur en charge}
				\label{graph:mee-charge-emitter}
			\end{figure}
		\end{column}

		\begin{column}{0.5\textwidth}
			\begin{figure}
				\includegraphics[width=.9\textwidth]{figure/tp34-single.png}
				\caption{Simulation de l'émetteur en charge}
				\label{graph:lts-charge-emitter}
			\end{figure}
		\end{column}


	\end{columns}
\end{frame}

\subsection{À la décharge}
\begin{frame}
	\frametitle{Mises en équation de la décharge}
\end{frame}

\begin{frame}
	\frametitle{Graphiques de décharge théoriques}
	\begin{columns}

		\begin{column}{0.5\textwidth}
			\begin{figure}
				
\includegraphics[width=.9\textwidth]{figure/todo.png}
				\caption{M.E.É. de l'émetteur en décharge}
				\label{graph:mee-decharge-emitter}
			\end{figure}
		\end{column}

		\begin{column}{0.5\textwidth}
			\begin{figure}
				
\includegraphics[width=.9\textwidth]{figure/todo.png}
				\caption{Simulation de l'émetteur en décharge}
				\label{graph:lts-decharge-emitter}
			\end{figure}
		\end{column}


	\end{columns}
\end{frame}

\subsection{Résultats expérimentaux}
\begin{frame}
	\frametitle{Résultats expérimentaux}
	\begin{columns}

		\begin{column}{0.5\textwidth}
			\begin{figure}
				
\includegraphics[width=.9\textwidth]{figure/todo.png}
				\caption{Mesure de l'émetteur en charge}
				\label{graph:osc-charge-emitter}
			\end{figure}
		\end{column}

		\begin{column}{0.5\textwidth}
			\begin{figure}
				
\includegraphics[width=.9\textwidth]{figure/todo.png}
				\caption{Mesure de l'émetteur en décharge}
				\label{graph:osc-decharge-emitter}
			\end{figure}
		\end{column}


	\end{columns}
\end{frame}
% vim: ft=tex
